% ============================================================
%  程式識讀 完整題庫（30 題）
%  依據《程式識讀題目範例 Python 題本》原文整理
% ============================================================

% 知識點標籤
\newcommand{\qskill}[1]{%
  \par\vspace{2pt}
  {\small\color{notebar}\textit{知識點：#1}}%
  \vspace{2pt}%
}

% 四個選項：每個都獨立一行，間距緊湊
\newcommand{\qchoices}[4]{%
  \vspace{3pt}
  \begin{enumerate}[label=(\Alph*), itemsep=1pt, topsep=1pt,
                    leftmargin=2.2em, parsep=0pt]
    \item #1
    \item #2
    \item #3
    \item #4
  \end{enumerate}
}

% -------- 題目 1 --------
\begin{qbox}{1}
給定右側函式 \texttt{f()}，當執行 \texttt{f(10)} 時，最終回傳結果為何？

\begin{minted}{python}
def f(i):
    if i > 0:
        if (i // 2) % 2 == 0:
            return f(i - 2) * i
        return f(i - 2) * (-i)
    else:
        return 1
\end{minted}

\qchoices{1}{3840}{-3840}{執行時導致無窮迴圈，不會停止執行}
\qskill{遞迴追蹤、整數除法 (//)、取餘 (\%)}
\end{qbox}

\medskip

% -------- 題目 2 --------
\begin{qbox}{2}
給定右側程式，若已知輸出的結果為 \texttt{[1][2][3][5][4][6]}，程式中的 \texttt{(?)} 應為下列何者？

\begin{minted}{python}
for i in range(5):
    j = 0
    while   (?):
        print(f'[{i + j}]', end = '')
        j = j + 2
\end{minted}

\qchoices{j<i}{j>i}{j<=i}{j>=i}
\qskill{while 條件填空、f-string 輸出}
\end{qbox}

\medskip

% -------- 題目 3 --------
\begin{qbox}{3}
給定右側函式 \texttt{f()}，已知 \texttt{f(14)}、\texttt{f(10)}、\texttt{f(6)} 分別回傳 25、18、10，函式中的 \texttt{(?)} 應為下列何者？

\begin{minted}{python}
def f(n):
    if n < 2:
        return n
    else:
        return n + f(  (?)  )
\end{minted}

\qchoices{(n+1)//2}{n//2}{(n-1)//2}{(n//2)+1}
\qskill{遞迴規律推導、整數除法}
\end{qbox}

\medskip

% -------- 題目 4 --------
\begin{qbox}{4}
給定右側程式片段，當程式執行完後，輸出結果為何？

\begin{minted}{python}
Q = [0] * 200
val = 0
count = 0
head = tail = 0
for i in range(1, 31):
    Q[tail] = i
    tail = tail + 1
while tail > head + 1:
    val = Q[head]
    head = head + 1
    count = count + 1
    if count == 3:
        count = 0
        Q[tail] = val
        tail = tail + 1
print(Q[head])
\end{minted}

\qchoices{9}{18}{27}{30}
\qskill{串列模擬、head/tail 指標（佇列概念）}
\end{qbox}

\medskip

% -------- 題目 5 --------
\begin{qbox}{5}
右側程式正確的輸出應該如下：

\begin{center}\ttfamily\small
\phantom{****}*\\
\phantom{**}*\ *\ *\\
*\ *\ *\ *\ *\\
*\ *\ *\ *\ *\ *\ *\\
*\ *\ *\ *\ *\ *\ *\ *\ *
\end{center}

在不修改右側程式之第 4 行及第 5 行程式碼的前提下，最少需修改幾行程式碼以得到正確輸出？

\begin{minted}[linenos=true]{python}
k = 4
m = 1
for i in range(5):
    print(' ' * k, end = '')
    print('*' * m)
    k = k - 1
    m = m + 1
\end{minted}

\qchoices{1}{2}{3}{4}
\qskill{迴圈結構、字串重複 (* 運算子)}
\end{qbox}

\medskip

% -------- 題目 6 --------
\begin{qbox}{6}
給定一整數列表 \texttt{a[0]、a[1]、...、a[99]} 且 \texttt{a[k]=3k+1}，以 \texttt{value=100} 呼叫以下兩函式，假設函式 \texttt{f1} 及 \texttt{f2} 之 \texttt{while} 迴圈主體分別執行 n1 與 n2 次（也就是說計算 \texttt{if} 敘述執行次數，不包含 \texttt{elif} 敘述），請問 n1 與 n2 之值為何？

\begin{minted}{python}
def f1(a, value):          def f2(a, value):
    r_value = -1               r_value = -1
    i = 0                      low = 0
    while i < 100:             high = 99
        n1 += 1                while low <= high:
        if a[i] == value:          n2 += 1
            r_value = i            mid = (low + high) // 2
            break                  if a[mid] == value:
        i = i + 1                      r_value = mid
    return r_value                     break
                                   elif a[mid] < value:
                                       low = mid + 1
                                   else:
                                       high = mid - 1
                               return r_value
\end{minted}

\qchoices{n1=33, n2=4}{n1=33, n2=5}{n1=34, n2=4}{n1=34, n2=5}
\qskill{線性搜尋 vs 二分搜尋、迴圈次數計算}
\end{qbox}

\medskip

% -------- 題目 7 --------
\begin{qbox}{7}
一個費式數列定義第一個數為 0，第二個數為 1，之後的每個數都等於前兩個數相加，如下所示：0、1、1、2、3、5、8、13、21、34、55、89……右列的程式用以計算並印出第 N 個（N≥2）費式數列的數值，請問 \texttt{(a)} 與 \texttt{(b)} 兩個空格的敘述（statement）應該為何？

\begin{minted}{python}
a = 0
b = 1
for i in range(2, N + 1):
    temp = b
    (a)
    a = temp
print( (b) )
\end{minted}

\begin{enumerate}[label=(\Alph*), itemsep=3pt, topsep=1pt, leftmargin=2.2em, parsep=0pt]
  \item \texttt{(a) f[i] = f[i-1]+f[i-2]} \quad \texttt{(b) f[N]}
  \item \texttt{(a) a = a+b} \quad \texttt{(b) a}
  \item \texttt{(a) b = a+b} \quad \texttt{(b) a}
  \item \texttt{(a) f[i] = f[i-1]+f[i-2]} \quad \texttt{(b) f[i]}
\end{enumerate}
\qskill{費式數列、暫存變數（temp）、迴圈實作}
\end{qbox}

\medskip

% -------- 題目 8 --------
\begin{qbox}{8}
給定右側函式 \texttt{f1()} 及 \texttt{f2()}。\texttt{f1(1)} 運算過程中，以下敘述何者為錯？

\begin{minted}{python}
def f1(m):                def f2(n):
    if m > 3:                 if n > 3:
        print(m)                  print(n)
        return                    return
    else:                     else:
        print(m)                  print(n)
        f2(m + 2)                 f1(n - 1)
        print(m)                  print(n)
\end{minted}

\qchoices{印出的數字最大的是 4}{\texttt{f1} 一共被呼叫二次}{\texttt{f2} 一共被呼叫三次}{數字 2 被印出兩次}
\qskill{互相遞迴、執行流程追蹤、呼叫次數}
\end{qbox}

\medskip

% -------- 題目 9 --------
\begin{qbox}{9}
右側程式片段擬以輾轉除法求 \texttt{i} 與 \texttt{j} 的最大公因數。請問 \texttt{while} 迴圈內容何者正確？（注意：\texttt{(low + high) // 2} 只取整數部分）

\begin{minted}{python}
i = 76
j = 48
while (i % j) != 0:
    ________
    ________
    ________
print(j)
\end{minted}

\begin{enumerate}[label=(\Alph*), itemsep=4pt, topsep=1pt, leftmargin=2.2em, parsep=0pt]
  \item \texttt{k = i \% j} \\ \texttt{i = j} \\ \texttt{j = k}
  \item \texttt{i = j} \\ \texttt{j = k} \\ \texttt{k = i \% j}
  \item \texttt{i = j} \\ \texttt{j = i \% k} \\ \texttt{k = i}
  \item \texttt{k = i} \\ \texttt{i = j} \\ \texttt{j = i \% k}
\end{enumerate}
\qskill{輾轉除法（GCD）、while 迴圈填空}
\end{qbox}

\medskip

% -------- 題目 10 --------
\begin{qbox}{10}
右側程式執行過後所輸出數值為何？

\begin{minted}{python}
count = 10
if count > 0:
    count = 11
if count > 10:
    count = 12
    if count % 3 == 4:
        count = 1
    else:
        count = 0
elif count > 11:
    count = 13
else:
    count = 14
if count:
    count = 15
else:
    count = 16
print(count)
\end{minted}

\qchoices{11}{13}{15}{16}
\qskill{獨立 if 追蹤、if/elif 執行順序、布林值判斷}
\end{qbox}

\medskip

% -------- 題目 11 --------
\begin{qbox}{11}
右側為一個計算 n 階乘的函式，請問該如何修改才會得到正確的結果？

\begin{minted}[linenos=true]{python}
def fun(n):
    fac = 1
    if n >= 0:
        fac = n * fun(n - 1)
    return fac
\end{minted}

\begin{enumerate}[label=(\Alph*), itemsep=1pt, topsep=1pt, leftmargin=2.2em, parsep=0pt]
  \item 第 2 行，改為 \texttt{fac = n}
  \item 第 3 行，改為 \texttt{if n > 0:}
  \item 第 4 行，改為 \texttt{fac = n * fun(n+1)}
  \item 第 4 行，改為 \texttt{fac = fac * fun(n-1)}
\end{enumerate}
\qskill{遞迴終止條件（base case）、階乘}
\end{qbox}

\medskip

% -------- 題目 12 --------
\begin{qbox}{12}
右側 \texttt{f()} 函式 \texttt{(a)}、\texttt{(b)}、\texttt{(c)} 處需分別填入哪些數字，方能使得 \texttt{f(4)} 輸出 \texttt{2468} 的結果？

\begin{minted}{python}
def f(n):
    p = 0
    i = n
    while i >= (a):
        p = 10 - (b) * i
        print(p, end='')
        i = i - (c)
\end{minted}

\qchoices{1, 2, 1}{0, 1, 2}{0, 2, 1}{1, 1, 1}
\qskill{while 條件與迴圈體填空、輸出預測}
\end{qbox}

\medskip

% -------- 題目 13 --------
\begin{qbox}{13}
右側 \texttt{g(4)} 函式呼叫執行後，回傳值為何？

\begin{minted}{python}
def f(n):
    if n > 3:
        return 1
    elif n == 2:
        return 3 + f(n + 1)
    else:
        return 1 + f(n + 1)

def g(n):
    j = 0
    for i in range(1, n):
        j = j + f(i)
    return j
\end{minted}

\qchoices{6}{11}{13}{14}
\qskill{遞迴 + for 迴圈混合追蹤、累加}
\end{qbox}

\medskip

% -------- 題目 14 --------
\begin{qbox}{14}
右側 \texttt{Mystery()} 函式 \texttt{else} 部分運算式應為何，才能使得 \texttt{Mystery(9)} 的回傳值為 34。

\begin{minted}{python}
def Mystery(x):
    if x <= 1:
        return x
    else:
        return ________
\end{minted}

\qchoices{x + Mystery(x-1)}{x * Mystery(x-1)}{Mystery(x-2) + Mystery(x+2)}{Mystery(x-2) + Mystery(x-1)}
\qskill{遞迴終止條件推導、逆向推算}
\end{qbox}

\medskip

% -------- 題目 15 --------
\begin{qbox}{15}
右側程式碼執行後，輸出結果為何？

\begin{minted}{python}
a = [1, 3, 5, 7, 9, 8, 6, 4, 2]
n = 9
for i in range(n):
    a[i], a[n-i-1] = a[n-i-1], a[i]
for i in range(n//2 + 1):
    print(a[i], a[n-i-1], end=' ')
\end{minted}

\qchoices{2468975319}{1357924689}{1234567899}{2468513799}
\qskill{串列同時賦值、反轉、索引計算}
\end{qbox}

\medskip

% -------- 題目 16 --------
\begin{qbox}{16}
右側函式以 \texttt{F(7)} 呼叫後回傳值為 12，則 \texttt{<condition>} 應為何？

\begin{minted}{python}
def F(a):
    if <condition>:
        return 1
    else:
        return F(a-2) + F(a-3)
\end{minted}

\qchoices{a < 3}{a < 2}{a < 1}{a < 0}
\qskill{遞迴終止條件填空、逆向推算}
\end{qbox}

\medskip

% -------- 題目 17 --------
\begin{qbox}{17}
右側是依據分數 \texttt{s} 評定等第的程式碼片段，正確的等第公式應為：90\textasciitilde100 判為 A 等、80\textasciitilde89 判為 B 等、70\textasciitilde79 判為 C 等、60\textasciitilde69 判為 D 等、0\textasciitilde59 判為 F 等。這段程式碼在處理 0\textasciitilde100 的分數時，有幾個分數的等第是錯的？

\begin{minted}{python}
if s >= 90:
    print('A')
elif s >= 80:
    print('B')
elif s > 60:
    print('D')
elif s > 70:
    print('C')
else:
    print('F')
\end{minted}

\qchoices{20}{11}{2}{10}
\qskill{if/elif 條件順序、邊界值（>  vs  >=）、錯誤計數}
\end{qbox}

\medskip

% -------- 題目 18 --------
\begin{qbox}{18}
右側程式片段執行後，\texttt{count} 的值為何？

\begin{minted}{python}
maze = [[1, 1, 1, 1, 1],
        [1, 0, 1, 0, 1],
        [1, 1, 0, 0, 1],
        [1, 0, 0, 1, 1],
        [1, 1, 1, 1, 1]]
count = 0
for i in range(1, 4):
    for j in range(1, 4):
        dir = [[-1,0], [0,1], [1,0], [0,-1]]
        for d in range(4):
            if maze[i+dir[d][0]][j+dir[d][1]] == 1:
                count = count + 1
\end{minted}

\qchoices{36}{20}{12}{3}
\qskill{二維串列、巢狀迴圈、方向向量}
\end{qbox}

\medskip

% -------- 題目 19 --------
\begin{qbox}{19}
右側程式片段中執行後若要印出下列圖案，\texttt{(a)} 的條件判斷式該如何設定？

\begin{center}\ttfamily\small
* * * * * *\\
* * * *\\
* *
\end{center}

\begin{minted}{python}
for i in range(4):
    for j in range(i):
        print(' ', end='')
    for k in range(6-2*i, (a), -1):
        print('*', end='')
    print('')
\end{minted}

\qchoices{2}{1}{0}{-1}
\qskill{range 參數填空、倒三角圖形}
\end{qbox}

\medskip

% -------- 題目 20 --------
\begin{qbox}{20}
下列程式碼是自動計算找零程式的一部分，程式碼中三個主要變數分別為 \texttt{Total}（購買總額）、\texttt{Paid}（實際支付金額）、\texttt{Change}（找零金額）。但是此程式片段有冗餘的程式碼，也就是移除該段程式碼之後不會影響程式的功能。請找出冗餘程式碼的區塊。

\begin{minted}{python}
Change = Paid - Total
print(f'500 : {(Change-Change%500)//500} pieces')
Change = Change % 500
print(f'100 : {(Change-Change%100)//100} coins')
Change = Change % 100
# A 區
print(f'50 : {(Change-Change%50)//50} coins')
Change = Change % 50
# B 區
print(f'10 : {(Change-Change%10)//10} coins')
Change = Change % 10
# C 區
print(f'5 : {(Change-Change%5)//5} coins')
Change = Change % 5
# D 區
print(f'1 : {(Change-Change%1)//1} coins')
Change = Change % 1
\end{minted}

\qchoices{冗餘程式碼在 A 區}{冗餘程式碼在 B 區}{冗餘程式碼在 C 區}{冗餘程式碼在 D 區}
\qskill{冗餘程式碼識別（\% 1 的性質）}
\end{qbox}

\medskip

% -------- 題目 21 --------
\begin{qbox}{21}
右側 \texttt{G()} 為遞迴函式，\texttt{G(3, 7)} 執行後回傳值為何？

\begin{minted}{python}
def G(a, x):
    if x == 0:
        return 1
    else:
        return a * G(a, x - 1)
\end{minted}

\qchoices{128}{2187}{6561}{1024}
\qskill{遞迴幂次、G(a, x) = a\textsuperscript{x}}
\end{qbox}

\medskip

% -------- 題目 22 --------
\begin{qbox}{22}
給定一個 1×8 的列表 \texttt{A}，\texttt{A = [0, 2, 4, 6, 8, 10, 12, 14]}。右側函式 \texttt{Search(x)} 真正目的是找到 \texttt{A} 之中大於 \texttt{x} 的最小值。然而，這個函式有誤。請問下列哪個函式呼叫可以測出函式有誤？

\begin{minted}{python}
A = [0, 2, 4, 6, 8, 10, 12, 14]
def Search(x):
    high = 7
    low = 0
    while high > low:
        mid = (high + low) // 2
        if A[mid] <= x:
            low = mid + 1
        else:
            high = mid
    return A[high]
\end{minted}

\qchoices{Search(-1)}{Search(0)}{Search(10)}{Search(16)}
\qskill{二分搜尋邏輯、邊界測試（錯誤觸發）}
\end{qbox}

\medskip

% -------- 題目 23 --------
\begin{qbox}{23}
右側函式兩個回傳式分別該如何撰寫，才能正確計算並回傳參數 \texttt{a}、\texttt{b} 之最大公因數（Greatest Common Divisor）？

\begin{minted}{python}
def GCD(a, b):
    r = a % b
    if r == 0:
        return ____
    return ____
\end{minted}

\qchoices{a,\ \ GCD(b,r)}{b,\ \ GCD(b,r)}{a,\ \ GCD(a,r)}{b,\ \ GCD(a,r)}
\qskill{GCD 遞迴填空、輾轉除法}
\end{qbox}

\medskip

% -------- 題目 24 --------
\begin{qbox}{24}
若 \texttt{A} 是一個可儲存 n 筆整數的列表，且資料儲存於 \texttt{A[0]}$\sim$\texttt{A[n-1]}。經過右側程式碼運算後，以下何者敘述不一定正確？

\begin{minted}{python}
A = [...]
p = q = A[0]
for i in range(1, n):
    if A[i] > p:
        p = A[i]
    if A[i] < q:
        q = A[i]
\end{minted}

\qchoices{p 是 A 列表資料中的最大值}{q 是 A 列表資料中的最小值}{q < p}{A[0] $\leq$ p}
\qskill{最大最小值追蹤、「不一定正確」邏輯陷阱}
\end{qbox}

\medskip

% -------- 題目 25 --------
\begin{qbox}{25}
右側 \texttt{F()} 函式執行時，若輸入依序為整數 0, 1, 2, 3, 4, 5, 6, 7, 8, 9，請問 \texttt{X} 列表的元素值依序為何？

\begin{minted}{python}
def F():
    X = [0] * 10
    for i in range(10):
        X[(i+2)%10] = int(input())
\end{minted}

\qchoices{0, 1, 2, 3, 4, 5, 6, 7, 8, 9}{2, 0, 2, 0, 2, 0, 2, 0, 2, 0}{9, 0, 1, 2, 3, 4, 5, 6, 7, 8}{8, 9, 0, 1, 2, 3, 4, 5, 6, 7}
\qskill{串列索引偏移 ((i+2)\%10)、模運算}
\end{qbox}

\medskip

% -------- 題目 26 --------
\begin{qbox}{26}
右側程式片段無法正確列印 20 次的 \texttt{"Hi!"}，請問下列哪一個修改方式仍無法正確列印 20 次的 \texttt{"Hi!"}？

\begin{minted}{python}
for i in range(0, 101, 5):
    print('Hi!')
\end{minted}

\begin{enumerate}[label=(\Alph*), itemsep=1pt, topsep=1pt, leftmargin=2.2em, parsep=0pt]
  \item 將 \texttt{range(0,101,5)} 改為 \texttt{range(0,20,1)}
  \item 將 \texttt{range(0,101,5)} 改為 \texttt{range(5,101,5)}
  \item 將 \texttt{range(0,101,5)} 改為 \texttt{range(0,100,5)}
  \item 將 \texttt{range(0,101,5)} 改為 \texttt{range(5,100,5)}
\end{enumerate}
\qskill{range 元素數量、off-by-one}
\end{qbox}

\medskip

% -------- 題目 27 --------
\begin{qbox}{27}
給定右側函式 \texttt{F()}，執行 \texttt{F()} 時哪一行程式碼永遠不會被執行到？

\begin{minted}{python}
def F(a):
    while a < 10:
        a = a + 5
    if a < 12:
        a = a + 2
    if a < 11:
        a = 5
\end{minted}

\qchoices{a = a + 5}{a = a + 2}{a = 5}{每一行都執行得到}
\qskill{while 後變數範圍、dead code 分析}
\end{qbox}

\medskip

% -------- 題目 28 --------
\begin{qbox}{28}
給定右側函式 \texttt{F()}，\texttt{F()} 執行完所回傳的 \texttt{x} 值為何？

\begin{minted}{python}
def F(n):
    x = 0
    for i in range(1, n+1):
        for j in range(i, n+1):
            k = 1
            while k <= n:
                x = x + 1
                k = k * 2
    return x
\end{minted}

\begin{enumerate}[label=(\Alph*), itemsep=2pt, topsep=1pt, leftmargin=2.2em, parsep=0pt]
  \item $n(n+1)\sqrt{\lfloor\log_2 n\rfloor}$
  \item $n^2(n+1)/2$
  \item $n(n+1)(\lfloor\log_2 n\rfloor + 1)/2$
  \item $n(n+1)/2$
\end{enumerate}
\qskill{三層巢狀迴圈、複雜度（log 因子）}
\end{qbox}

\medskip

% -------- 題目 29 --------
\begin{qbox}{29}
右側程式擬找出列表 \texttt{A} 中的最大值（M）和最小值（N）。不過，這段程式碼有誤，請問 \texttt{A} 初始值如何設定就可以測出程式有誤？

\begin{minted}{python}
M, N, s = -1, 101, 3
A = [80, 90, 100]
for i in range(s):
    if A[i] > M:
        M = A[i]
    elif A[i] < N:
        N = A[i]
print(f'M = {M}, N = {N}')
\end{minted}

\qchoices{\text{[90, 80, 100]}}{\text{[80, 90, 100]}}{\text{[100, 90, 80]}}{\text{[90, 100, 80]}}
\qskill{elif 邏輯漏洞、最大最小值、錯誤測資設計}
\end{qbox}

\medskip

% -------- 題目 30 --------
\begin{qbox}{30}
經過運算後，右側程式的輸出為何？

\begin{minted}{python}
a = [0] * 101
b = [i for i in range(101)]
for i in range(1, 101):
    a[i] = b[i] + a[i - 1]
print(a[50] - a[30])
\end{minted}

\qchoices{1275}{20}{1000}{810}
\qskill{list comprehension、前綴和（prefix sum）}
\end{qbox}

\bigskip

% -------- 答案速查表 --------
\begin{center}
\small
\begin{tabular}{ccccccccccc}
\toprule
題號 & 1 & 2 & 3 & 4 & 5 & 6 & 7 & 8 & 9 & 10 \\
答案 & C & A & B & B & A & D & C & C & A & D \\
\midrule
題號 & 11 & 12 & 13 & 14 & 15 & 16 & 17 & 18 & 19 & 20 \\
答案 & B & A & C & D & C & D & B & B & C & D \\
\midrule
題號 & 21 & 22 & 23 & 24 & 25 & 26 & 27 & 28 & 29 & 30 \\
答案 & B & D & B & C & D & D & C & C & B & D \\
\bottomrule
\end{tabular}
\end{center}
